\documentclass[11pt]{article}
\usepackage[a4paper,margin=1in]{geometry}
\usepackage{amsmath,amssymb,amsthm}
\usepackage{array}
\usepackage{hyperref}

\newcommand{\dd}{\mathrm{d}}
\newcommand{\E}{\mathbb{E}}
\newcommand{\sgn}{\operatorname{sgn}}
\newcommand{\He}{\mathrm{He}}
\newcommand{\vt}{v_t}
\newcommand{\avg}[1]{\langle #1\rangle}
\renewcommand{\Re}{\operatorname{Re}}
\renewcommand{\Im}{\operatorname{Im}}

\newtheorem{theorem}{Theorem}
\newtheorem{lemma}[theorem]{Lemma}
\newtheorem{corollary}[theorem]{Corollary}
\theoremstyle{remark}
\newtheorem{remark}[theorem]{Remark}

\title{What is missing in plasma closure theory,\\
and an exact closure at every collisionality for a model case}
\author{Working note}
\date{16 November 2026}

\begin{document}
\maketitle

\noindent\emph{Summary.} Part~I lists the gaps in the theory of plasma
closures, in the field and in this project's own notes. Part~II addresses
the gap singled out as most relevant to the benchmark, the inhomogeneous
case at intermediate collisionality, and proves a theorem for the simplest
model in which it can be posed: the linearized one-dimensional BGK equation.
The heat-flux closure is obtained exactly at every wavenumber and collision
frequency, it is shown to be dissipative and always below Fourier's law, and
both limits are recovered with explicit error bounds, one of them sharp. The
theorem is linear, uses a velocity-independent collision frequency, and is
for a potential-driven response; Section~\ref{sec:limits} says what that
leaves out, which includes the mechanism behind the early failure of
Spitzer--H\"arm transport.

\tableofcontents

\part{What is missing}

\section{In the field}

\begin{itemize}
  \item \textbf{Particles to Vlasov with Coulomb forces in three
    dimensions.} Available with an $N$-dependent cutoff
    \cite{LazaroviciPickl}; a 2026 preprint claims two dimensions without
    cutoff \cite{DuerinckxJabin}.
  \item \textbf{Particles to Landau or Lenard--Balescu on the relaxation
    time.} For the true dynamics the best result is linearized, at
    equilibrium, and valid only up to a small power of $N$
    \cite{DuerinckxSR}. Time $N$ is reached for a toy model with an added
    velocity diffusion \cite{DuerinckxLeBihan}.
  \item \textbf{A chained result from particles to fluid.} The hard-sphere
    chain of Deng, Hani and Ma \cite{DHM2} has no plasma counterpart; plasma
    fluid limits start from a kinetic equation \cite{GuoJang}.
  \item \textbf{Quasilinear and weak-turbulence theory for the
    self-consistent problem.} Proofs exist for a prescribed field only
    \cite{BardosBesse}.
  \item \textbf{A theorem behind Landau-fluid closures.} The
    Hammett--Perkins closure \cite{HP90} is a matched linear response, with
    no error bound in the nonlinear problem and no treatment of bounded
    field lines.
  \item \textbf{Intermediate collisionality.} Nonlocal heat-flux models are
    constructed to interpolate between the collisional and collisionless
    regimes; for the physical collision operator nothing proves that a given
    model recovers both limits with controlled error.
  \item \textbf{Strong coupling.} Dusty plasmas lie outside every expansion
    above.
\end{itemize}

\section{In this project's strategy for the collisional kinetic limit}

Refer to \texttt{hani-to-algorithms/docs/plasma\_extension.tex}.
\begin{itemize}
  \item \textbf{The bootstrap is open.} It must show that phase-mixed
    information cannot return, using diffusion that the fluctuations generate
    themselves. Shot noise sits exactly at the reported low-regularity echo
    threshold, so the margin has to come from independence of the particles,
    and there is no argument for it yet.
  \item \textbf{The one-layer estimate is not done.} It needs mean-square
    bounds for a diagram expansion with non-Gaussian data on times up to
    $N^{1/3}$.
  \item \textbf{The proved theorem is a model.} Orbit decorrelation is
    proved for a prescribed Gaussian Markovian field, not for the field
    generated by the particles.
  \item \textbf{Stability of the limit equation} within the Penrose-stable
    class is assumed.
  \item \textbf{Inhomogeneous states} are not treated.
  \item \textbf{The Coulomb split} at the interparticle spacing rests on an
    unchecked estimate.
\end{itemize}

\section{As evidence}

\begin{itemize}
  \item No proof in these notes has been checked by a second reader.
  \item The echo experiment gave no result.
  \item Nothing connects the kinetic limit to a fluid closure with a
    quantified defect. Part~II is a first step for one model.
  \item Algorithm~D is not implemented.
\end{itemize}

\part{An exact closure at every collisionality}

\section{Setting}

Units: $n_0 = m = \vt = 1$ with $\vt = \sqrt{T_0/m}$. Let
$M(v) = (2\pi)^{-1/2}e^{-v^2/2}$ and $\avg{\varphi} = \int\varphi M\,\dd v$.
Write the distribution as $M(1+h)$. The linearized BGK equation with
collision frequency $\nu\ge0$ and acceleration $a(t,x)$ is
\begin{equation}
  \partial_t h + v\,\partial_x h = \nu\,(\Pi h - h) + a\,v,
  \qquad h|_{t=0} = 0,
  \label{eq:bgk}
\end{equation}
on $x\in\mathbb{T}$ or $\mathbb{R}$, where $\Pi$ is the orthogonal projection
in $L^2(M\,\dd v)$ onto $\operatorname{span}\{1,v,v^2\}$:
\[
  \Pi h = n + u\,v + \tfrac12\,T\,(v^2-1),\qquad
  n = \avg{h},\quad u = \avg{vh},\quad T = \avg{(v^2-1)h}.
\]
These are the density, velocity and temperature perturbations, and the heat
flux is $q = \avg{(v^3-3v)h}$. The operator conserves mass, momentum and
energy. Since $\Pi-1$ is non-positive,
$\|h(t)\|\le\int_0^t\|a(\sigma)\|\,\dd\sigma$ in $L^2(\dd x\,M\dd v)$, so for
$a\in L^\infty(0,\infty;L^2_x)$ the Fourier--Laplace transform
\[
  \hat h(k,\omega,v) = \int_0^\infty e^{i\omega t}\int e^{-ikx}h(t,x,v)\,\dd x\,\dd t
\]
exists for $\Im\omega>0$, and
\begin{equation}
  (\nu - i\omega + ikv)\,\hat h = \nu\,\Pi\hat h + \hat a\,v .
  \label{eq:fl}
\end{equation}

For $z\in\mathbb{C}\setminus\mathbb{R}$ define
\begin{equation}
  W(z) = \Bigl\langle\frac{1}{v-z}\Bigr\rangle,\qquad
  R(z) = 1 + zW(z),\qquad
  \mathcal{E}(z) = 1 - R(z) + z^2R(z).
\end{equation}
$W(z) = Z(z/\sqrt2)/\sqrt2$ with $Z$ the plasma dispersion function, and $R$
is the collisionless response function. Define
\begin{equation}
  \Phi(z) = \frac{z\,\bigl(1 - 3R(z) + z^2R(z)\bigr)}{\mathcal{E}(z)} .
  \label{eq:Phi}
\end{equation}
Throughout, $t$ denotes $v^2$ under the Gaussian law (so $\E t = 1$,
$\E t^2 = 3$, $\E t^3 = 15$, $\E t^4 = 105$), $s = y^2$, and
$G(s) = \E[1/(s+t)]$.

\section{Results}

\begin{theorem}[Exact linear closure]\label{thm:exact}
Let $\nu\ge0$, $k\ne0$ and $\Im\omega>0$. Every solution of \eqref{eq:bgk}
satisfies
\begin{equation}
  \hat q(k,\omega) = \sgn(k)\;\Phi\!\Bigl(\frac{\omega+i\nu}{|k|}\Bigr)\,
  \hat T(k,\omega) .
  \label{eq:exact}
\end{equation}
$\Phi$ is analytic in the upper half-plane. The collisional closure is the
collisionless one evaluated at the complex phase velocity
$(\omega+i\nu)/|k|$.
\end{theorem}

\begin{theorem}[Quasi-static closure at every collisionality]\label{thm:theta}
For $y>0$, $\Phi(iy) = -i\,\theta(y)$ with
\begin{equation}
  \theta(y) = \frac{1+y^2}{y\,(3+y^2)}\;
  \frac{\bigl\langle \He_3(v)^2/(v^2+y^2)\bigr\rangle}
       {\bigl\langle \He_2(v)^2/(v^2+y^2)\bigr\rangle},
  \qquad \He_2 = v^2-1,\quad \He_3 = v^3-3v .
  \label{eq:theta}
\end{equation}
Consequently, in the limit $\omega\to0$ of \eqref{eq:exact},
$\hat q = -i\,\sgn(k)\,\theta(\nu/|k|)\,\hat T$, and:
\begin{enumerate}
  \item[(a)] \emph{Positivity.} $\theta(y)>0$ for all $y>0$.
  \item[(b)] \emph{Below Fourier's law.} $\theta(y) < 3/y$ for all $y>0$.
  \item[(c)] \emph{Collisional limit, sharp.}
    $0 < 3/y - \theta(y) \le 12/y^3$, and $12$ cannot be improved:
    $y^3(3/y-\theta(y))\to12$ as $y\to\infty$.
  \item[(d)] \emph{Collisionless limit.}
    $|\theta(y) - \sqrt{8/\pi}| \le 5y$ for $0<y\le\tfrac14$.
\end{enumerate}
\end{theorem}

\begin{corollary}[Physical units]\label{cor:units}
Let $\lambda = \vt/\nu$ be the mean free path. The quasi-static closure is
$\hat q_k = -\chi(k)\,ik\,\hat T_k$ with
$\chi(k) = n_0\vt\,\theta\bigl(1/(|k|\lambda)\bigr)/|k|$, and
\begin{align*}
  &0 < \chi(k) < \kappa_{\mathrm F}, \qquad
   \kappa_{\mathrm F} = \frac{3n_0T_0}{m\nu}, \\
  &0 < 1 - \frac{\chi(k)}{\kappa_{\mathrm F}} \le 4\,(k\lambda)^2, \\
  &\Bigl|\,\chi(k) - \sqrt{\tfrac{8}{\pi}}\,\frac{n_0\vt}{|k|}\Bigr|
   \le \frac{5\,n_0\vt}{|k|}\cdot\frac{1}{|k|\lambda}
   \qquad (|k|\lambda\ge4).
\end{align*}
The first line says the exact flux never exceeds the Fourier-law flux. The
second gives the relative error of the Navier--Stokes--Fourier closure. The
third recovers the Hammett--Perkins closure.
\end{corollary}

\begin{corollary}[Dissipativity]\label{cor:diss}
For the quasi-static closure on $\mathbb{T} = [0,2\pi)$, with
$\hat T_k = \int e^{-ikx}T\,\dd x$,
\[
  \int q\,\partial_xT\,\dd x
  = -\frac{1}{2\pi}\sum_{k\ne0}|k|\,\theta(\nu/|k|)\,|\hat T_k|^2 \le 0
\]
for every temperature profile, with equality only for constants.
\end{corollary}

\section{Proofs}

\subsection{Three identities}

\begin{lemma}\label{lem:id}
Let $s>0$, $P_2(t) = t^2-6t+3$ and $p_0 = s^2+6s+3$. Put
$N = 2+s-s(3+s)G$ and $D = s[(1+s)G-1]$. Then
\begin{align}
  D &= \frac{s}{1+s}\;\E\Bigl[\frac{(t-1)^2}{s+t}\Bigr], \label{eq:idD}\\
  N &= \frac{1}{3+s}\;\E\Bigl[\frac{t\,(t-3)^2}{s+t}\Bigr], \label{eq:idN}\\
  3D - sN &= \frac{s}{p_0}\;\E\Bigl[\frac{P_2(t)^2}{s+t}\Bigr]. \label{eq:idP}
\end{align}
In particular $D>0$, $N>0$ and $3D-sN>0$.
\end{lemma}

\begin{proof}
\emph{\eqref{eq:idD}.} $(t-1)^2 = (t+s)^2 - 2(1+s)(t+s) + (1+s)^2$, so
$\E[(t-1)^2/(s+t)] = (1+s) - 2(1+s) + (1+s)^2G = (1+s)[(1+s)G-1]$.

\emph{\eqref{eq:idN}.}
$t(t-3)^2 = (t+s)\bigl(t^2-(6+s)t+(3+s)^2\bigr) - s(3+s)^2$, and
$\E[t^2-(6+s)t+(3+s)^2] = 3-(6+s)+(3+s)^2 = (s+2)(s+3)$. Hence
$\E[t(t-3)^2/(s+t)] = (s+2)(s+3) - s(3+s)^2G = (3+s)N$.

\emph{\eqref{eq:idP}.} $P_2$ is orthogonal to polynomials of degree at most
one: $\E P_2 = 3-6+3 = 0$ and $\E[tP_2] = 15-18+3 = 0$. Write
$P_2(t) = p_0 + (t+s)\,r(t)$ with $r(t) = t-s-6$, noting $P_2(-s) = p_0$.
Then
\[
  \E\Bigl[\frac{P_2^2}{s+t}\Bigr]
  = \E\Bigl[P_2\Bigl(\frac{p_0}{s+t} + r\Bigr)\Bigr]
  = p_0\,\E\Bigl[\frac{P_2}{s+t}\Bigr]
  = p_0\bigl(p_0G + \E r\bigr) = p_0\bigl(p_0G - (s+5)\bigr).
\]
On the other hand
$3D-sN = sG\,[3(1+s)+s(3+s)] - s(5+s) = s\,[p_0G-(s+5)]$.
\end{proof}

\subsection{Proof of Theorem~\ref{thm:exact}}

\emph{Conservation laws.} Taking the moments of \eqref{eq:fl} against $1$ and
$v^2-1$, and using $\avg{(v^2-1)v\,\hat h} = \hat q + 2\hat u$,
\begin{equation}
  -i\omega\,\hat n + ik\,\hat u = 0, \qquad
  -i\omega\,\hat T + 2ik\,\hat u + ik\,\hat q = 0 .
  \label{eq:cons}
\end{equation}
With $c = \omega/k$ these give $\hat u = c\,\hat n$ and
$\hat q = c\,\hat T - 2c\,\hat n$.

\emph{Kinetic relation, $k>0$.} Put $z = (\omega+i\nu)/k$, so $\Im z>0$ and
$\nu-i\omega+ikv = ik(v-z)$. From \eqref{eq:fl},
\[
  \hat h = \frac{A + Bv + C(v^2-1)}{ik\,(v-z)},\qquad
  A = \nu\hat n,\quad B = \nu\hat u + \hat a,\quad C = \tfrac12\nu\hat T .
\]
Let $I_m = \avg{v^m/(v-z)}$. From $v^m/(v-z) = v^{m-1} + z\,v^{m-1}/(v-z)$,
\[
  I_0 = W,\quad I_1 = R,\quad I_2 = zR,\quad I_3 = 1+z^2R,\quad
  I_4 = z+z^3R .
\]
The moments of $\hat h$ against $1$ and $v^2$ are
\begin{align*}
  ik\,\hat n &= A\,W + B\,R + C\,(zR-W),\\
  ik\,(\hat n+\hat T) &= A\,zR + B\,(1+z^2R) + C\,(z+z^3R-zR).
\end{align*}
Multiply the first by $1+z^2R$, the second by $R$, and subtract, which
eliminates $B$ and with it the force:
\[
  ik\bigl(\hat n\,\mathcal{E} - \hat T R\bigr)
  = A\bigl[W(1+z^2R)-zR^2\bigr]
  + C\bigl[(zR-W)(1+z^2R) - R(z+z^3R-zR)\bigr].
\]
Using $zW = R-1$ repeatedly, the first bracket is $W - zR$ and the second is
$zR - W$; and $z(zR-W) = \mathcal{E}$. Therefore
\[
  ikz\bigl(\hat n\,\mathcal{E} - \hat TR\bigr)
  = \nu\,\mathcal{E}\bigl(\tfrac12\hat T - \hat n\bigr),
  \qquad\text{that is}\qquad
  i\omega\,\mathcal{E}\,\hat n = \bigl(ikzR + \tfrac12\nu\mathcal{E}\bigr)\hat T,
\]
because $ikz+\nu = i\omega$.

\emph{Combination.} Since $2c = (2/ik)\,i\omega$,
\[
  \mathcal{E}\,\hat q = c\,\mathcal{E}\,\hat T - \frac{2}{ik}\,i\omega\,\mathcal{E}\,\hat n
  = \Bigl[c\,\mathcal{E} - 2zR + \frac{i\nu}{k}\,\mathcal{E}\Bigr]\hat T
  = z\,(\mathcal{E}-2R)\,\hat T ,
\]
using $c + i\nu/k = z$. No division by $\omega$ or $\hat a$ has occurred.

\emph{$\mathcal{E}$ does not vanish.} By evenness of $M$,
$\avg{p(v^2)/(v-z)} = z\,\avg{p(v^2)/(v^2-z^2)}$ for polynomials $p$. With
$\xi = z^2\in\mathbb{C}\setminus[0,\infty)$ this gives
$R = \E[t/(t-\xi)]$ and
\[
  \mathcal{E}(z) = \xi\;\E\Bigl[\frac{t-1}{t-\xi}\Bigr].
\]
If $\xi = \alpha+i\beta$ with $\beta\ne0$ and $\E[(t-1)/(t-\xi)] = 0$, the
imaginary and real parts give $\E[(t-1)/|t-\xi|^2] = 0$ and
$\E[(t-1)(t-\alpha)/|t-\xi|^2] = 0$; subtracting $(1-\alpha)$ times the first
from the second yields $\E[(t-1)^2/|t-\xi|^2] = 0$, which is impossible. If
$\xi = -s<0$, then
$\mathcal{E} = -s\,\E[(t-1)/(t+s)] = s\,[(1+s)G - 1] = D$, which is positive
by \eqref{eq:idD}. Hence
$\mathcal{E}\ne0$ in the upper half-plane, $\Phi$ is analytic there, and
$\hat q = \Phi(z)\hat T$.

\emph{$k<0$.} Now $z = (\omega+i\nu)/k$ has $\Im z<0$; the algebra above is
unchanged. Substituting $v\to-v$ gives $W(z) = -W(-z)$, hence
$R(z) = R(-z)$, $\mathcal{E}(z) = \mathcal{E}(-z)$ and
$\Phi(z) = -\Phi(-z)$ with $-z = (\omega+i\nu)/|k|$. \qed

\subsection{Proof of Theorem~\ref{thm:theta}}

For $z = iy$, $\xi = -s$, so $R = \E[t/(t+s)] = 1-sG$ is real. Then
$\mathcal{E}(iy) = 1-(1+s)(1-sG) = D$ and
$1 - 3R + z^2R = 1-(3+s)(1-sG) = -N$. Thus
$\Phi(iy) = -i\,yN/D$, and $\theta = yN/D$. Formula \eqref{eq:theta} follows
from \eqref{eq:idD} and \eqref{eq:idN}, since $t(t-3)^2 = \He_3(v)^2$ and
$(t-1)^2 = \He_2(v)^2$.

(a) Both averages in \eqref{eq:theta} are positive.

(b) $3/y - \theta = (3D-sN)/(yD) > 0$ by Lemma~\ref{lem:id}.

(c) Since $1/(s+t)\le1/s$ and
$\E P_2^2 = 105-180+126-36+9 = 24$, \eqref{eq:idP} gives
$3D-sN\le24/p_0$. Since $3D-sN>0$, \eqref{eq:idP} also gives
$G > (s+5)/p_0$, hence
$D = s[(1+s)G-1] > s\,[(1+s)(s+5)-p_0]/p_0 = 2s/p_0$. Therefore
\[
  \frac3y - \theta = \frac{3D-sN}{yD} \le \frac{24/p_0}{y\cdot2s/p_0}
  = \frac{12}{y^3}.
\]
For sharpness, expand $G = s^{-1}(1 - s^{-1} + 3s^{-2} - 15s^{-3} + 105s^{-4} - \dots)$:
$N = 6s^{-2} - 60s^{-3} + O(s^{-4})$ and $D = 2s^{-1} - 12s^{-2} + O(s^{-3})$,
so $\theta = (3/y)(1 - 4/y^2 + O(y^{-4}))$.

(d) Let $w(y) = y\,G(y^2)$. Since
$y/(v^2+y^2) = \int_0^\infty e^{-y\tau}\cos(v\tau)\,\dd\tau$ and
$\avg{\cos(v\tau)} = e^{-\tau^2/2}$,
\[
  w(y) = \int_0^\infty e^{-y\tau-\tau^2/2}\,\dd\tau .
\]
So $w$ is decreasing, $w(0) = w_0 := \sqrt{\pi/2}$, and
$w'(y)\ge-\int_0^\infty\tau e^{-\tau^2/2}\dd\tau = -1$; hence
$w_0 - y\le w(y)\le w_0$. In terms of $w$,
\[
  \theta = \frac{2+s-(3+s)\,yw}{(1+s)\,w - y} .
\]
For $0<y\le\tfrac14$ the denominator is at least $w_0-2y>0$ and at most
$(1+s)w_0$, and the numerator lies between $2-(3+s)yw_0$ and $2+s$. Hence
\begin{align*}
  \theta - \frac{2}{w_0} &\le \frac{2+s}{w_0-2y} - \frac{2}{w_0}
    = \frac{y\,(yw_0+4)}{w_0(w_0-2y)} \le 4.6\,y,\\
  \frac{2}{w_0} - \theta &\le \frac{2}{w_0} - \frac{2-(3+s)yw_0}{(1+s)w_0}
    = \frac{y\,[2y+(3+s)w_0]}{(1+s)w_0} \le 3.5\,y ,
\end{align*}
using $w_0 = 1.2533\ldots$ and $y\le\tfrac14$. Finally
$2/w_0 = \sqrt{8/\pi}$. \qed

\subsection{Proofs of the corollaries}

Corollary~\ref{cor:units}: restore units with $y = \nu/(|k|\vt) = 1/(|k|\lambda)$;
then $\chi/\kappa_{\mathrm F} = y\theta(y)/3$ and (b), (c), (d) translate
directly. Corollary~\ref{cor:diss}: Parseval and (a).

\section{Numerical checks}

Scratch runs on 16 November 2026, not part of the repository.
\begin{itemize}
  \item \textbf{Independent solve.} The linearized BGK equation for one
    Fourier mode with harmonic forcing was solved in a Hermite basis
    truncated at 3000 modes, and $\hat q/\hat T$ compared with
    \eqref{eq:exact}. For $(k,\omega,\nu) = (1,0.3,2)$, $(1,0.7,0.5)$,
    $(-1,0.7,0.5)$, $(0.5,1.2,1)$ the difference was below $10^{-12}$; for
    $(2,0.1,0.3)$ it was $4\times10^{-7}$ and for $(-3,0.4,0.2)$ it was
    $3\times10^{-3}$. These two are at the weakest collisionality, where a
    Hermite truncation converges slowly; with $48000$ modes both differences
    fall below $10^{-12}$, so the discrepancy is truncation error.
  \item \textbf{Bounds.} On 281 logarithmically spaced points in
    $10^{-3}\le y\le10^4$, evaluated from \eqref{eq:theta} by quadrature,
    (a)--(d) hold, and $y^3(3/y-\theta)/12$ rises to $0.99999$ at $y=10^3$.
    An earlier evaluation from the equivalent formula $yN/D$ in double
    precision appeared to violate (a)--(c) for $y\gtrsim10^2$; that was
    cancellation error in $N$ and $D$, and it is the reason
    \eqref{eq:theta} is stated in a cancellation-free form.
\end{itemize}

\paragraph{Observed but not proved.} On the same grid, $\theta$ is
decreasing, $\theta(y)<\sqrt{8/\pi}$, and the harmonic-mean interpolant
\[
  \theta_{\mathrm H}(y) = \Bigl(\sqrt{\pi/8} + y/3\Bigr)^{-1}
\]
satisfies $1.000\le\theta/\theta_{\mathrm H}\le1.232$. If true in general,
the first two say the exact flux is below both the Fourier and the
collisionless values, and the third says a flux-limiter-type interpolation
underestimates the exact conductivity by at most about 19\%.

\begin{center}
\begin{tabular}{rrrr}
  $y = 1/(|k|\lambda)$ & $\theta(y)$ & $\sqrt{8/\pi}$ & $3/y$ \\ \hline
  0.01 & 1.5912 & 1.5958 & 300 \\
  0.1 & 1.5512 & 1.5958 & 30 \\
  1 & 1.2117 & 1.5958 & 3 \\
  3 & 0.7570 & 1.5958 & 1 \\
  10 & 0.2890 & 1.5958 & 0.3 \\
  100 & 0.029988 & 1.5958 & 0.03 \\
\end{tabular}
\end{center}

\section{Prior art}
\label{sec:prior}

Checked by a literature search on 16 November 2026. Most papers were
consulted at the level of abstract or reference list, so ``not found'' is
weaker than ``absent''.
\begin{itemize}
  \item \textbf{The shifted argument is standard.} Kinetic responses for a
    conserving BGK operator are written with
    $\zeta = (\omega+i\nu)/(\sqrt2|k|\vt)$, for example by Sharma, Hammett
    and Quataert \cite{SHQ03}. Their individual responses carry additional
    $\nu$-dependent terms from the conserving part. The statement that these
    terms cancel in the ratio $\hat q/\hat T$, leaving the collisionless
    ratio at the shifted argument, was not found.
  \item \textbf{Collisional Landau-fluid closures are interpolations.}
    Snyder, Hammett and Dorland remark that one could rederive the kinetic
    response with collisions and match it, and adopt a simpler procedure
    instead \cite{SHD97}; later closures are Pad\'e or harmonic-type fits
    \cite{JosephDimits,JiHeld}. A closed-form quasi-static coefficient at
    arbitrary collisionality for BGK was not found.
  \item \textbf{Exact hydrodynamics for linear BGK.} Kogelbauer and Karlin
    construct the slow invariant manifold of the three-dimensional linear
    BGK equation and show that each hydrodynamic branch exists only below a
    critical wavenumber of order $1/\lambda$ \cite{KK}. That object is an
    eigen-closure and ceases to exist at intermediate collisionality; the
    transfer function \eqref{eq:exact} is a different object, defined for
    all $k$, at the price of depending on the drive
    (Section~\ref{sec:limits}).
  \item \textbf{Bounds.} The inequality $\chi(k)<\kappa_{\mathrm F}$, the
    error bound with constant $4$, and the representation \eqref{eq:theta}
    were not found. The nearest statement is asymptotic and for the Landau
    operator: $Q\approx Q_{\mathrm{SH}}(1 - b\,Zk^2\lambda_{ei}^2)$ with
    $b\approx264$ in the Lorentz limit, due to Bychenkov and coworkers as
    quoted in \cite{Brodrick}.
  \item \textbf{Harmonic-mean interpolation.} Forms
    $\kappa_{\mathrm{SH}}/(1+a\,k\lambda)$ are standard. For the Landau
    operator they are reported to \emph{overestimate} the kinetic
    conductivity by more than 25\% near $k\lambda\approx0.5$
    \cite{Brodrick}. In the BGK model the harmonic mean is, numerically, a
    lower bound. The sign of the error therefore depends on the collision
    operator.
  \item \textbf{Not searched.} The gas-kinetic literature of the 1960s on
    sound propagation and density correlations in BGK models uses the same
    algebra and is the most likely place for an earlier appearance of
    Theorem~\ref{thm:exact}.
\end{itemize}

\begin{remark}
The collisionless limit here is the Hammett--Perkins value. A different
quasi-static collisionless coefficient, 40\% smaller, is attributed to the
closure of Chang and Callen in \cite{Brodrick}; that closure is for a
three-dimensional velocity space with a different collision model, and the
two are not expected to agree.
\end{remark}

\section{What the theorem does not cover}
\label{sec:limits}

\begin{enumerate}
  \item \textbf{Velocity-dependent collisions.} Here $\nu$ is constant, and
    Fourier's law is accurate to $4(k\lambda)^2$: it fails only when
    $k\lambda$ is of order one. In a real plasma the mean free path grows as
    $v^4$ and the heat-carrying electrons are several times faster than
    thermal, so the local closure fails at much smaller $k\lambda$. That
    mechanism is absent from this model. Extending the theorem to
    $\nu(v)$, or to the Lenard--Bernstein operator used as ground truth in
    the tournament note, is the natural next step and is not done here.

    A rough estimate shows the size of the effect. Iterating
    $h\mapsto-(v/\nu)\,\partial_xh$ on the temperature gradient and ignoring
    the conservation corrections gives a relative correction to Fourier's
    law of $-(k/\nu)^2\,\avg{\He_3v^3\He_2}/\avg{\He_3v\He_2} = -9\,(k/\nu)^2$
    for constant $\nu$, to be compared with the exact $-4\,(k/\nu)^2$. The
    same estimate with $\nu(v) = \nu_0|v|^{-3}$ replaces each factor $v$ by
    $v|v|^3$ and gives
    \[
      \frac{\avg{|v|^{17}-4|v|^{15}+3|v|^{13}}}
           {\avg{|v|^{9}-4|v|^{7}+3|v|^{5}}}\,
      \Bigl(\frac{k}{\nu_0}\Bigr)^2
      \approx 3.6\times10^4\,\Bigl(\frac{k}{\nu_0}\Bigr)^2 ,
    \]
    which reaches order one at $k/\nu_0\approx5\times10^{-3}$, where
    $1/\nu_0$ is the mean free path at the thermal speed. This is a
    heuristic, not a theorem, and it is good to a factor of a few at best;
    but it puts the breakdown of the local closure two orders of magnitude
    earlier in $k\lambda$ than in the constant-$\nu$ model. The
    nonlocal-transport literature attributes the failure of Spitzer--H\"arm
    at Knudsen numbers near $10^{-2}$ to exactly this: heat is carried by
    electrons near four times the thermal speed, whose mean free path is
    more than a hundred times the thermal one \cite{Brodrick}.
  \item \textbf{Linearity.} The result is a linear response about a uniform
    Maxwellian. Steep gradients with large temperature contrast are outside
    it.
  \item \textbf{The closure is not unique.} \eqref{eq:exact} holds for a
    response driven by a force from rest. Other drives, or non-zero initial
    data, add a term not proportional to $\hat T$. The quasi-static closure
    is the $\omega\to0$ value of a frequency-dependent relation; its defect
    in a time-dependent problem is $\Phi((\omega+i\nu)/|k|)-\Phi(i\nu/|k|)$,
    which is not bounded here.
  \item \textbf{One dimension, one species, no magnetic field.}
  \item \textbf{Real-space kernel.} The closure is
    $q = -\int K(x-x')\,\partial_{x'}T\,\dd x'$ with
    $\hat K(k) = \chi(k) > 0$. Positivity of $K$ itself, and its behaviour
    on a bounded domain, are not addressed.
\end{enumerate}

\section{Use in the benchmark}

For the linearized BGK model, \eqref{eq:exact} and \eqref{eq:theta} are an
exact reference at every collisionality. They give three tests for any
closure, classical or learned: recovery of $\kappa_{\mathrm F}$ with error
$\le4(k\lambda)^2$; recovery of the Hammett--Perkins coefficient as
$k\lambda\to\infty$; and the inequality $\chi(k)<\kappa_{\mathrm F}$ at every
$k$. A learned closure that violates the inequality in this model is wrong
there regardless of its pointwise error elsewhere.

\begin{thebibliography}{99}
\small
\bibitem{HP90} G.~W. Hammett, F.~W. Perkins, \emph{Fluid moment models for
Landau damping with application to the ion-temperature-gradient
instability}, Phys.\ Rev.\ Lett.\ \textbf{64} (1990) 3019.
\bibitem{DHM2} Y.~Deng, Z.~Hani, X.~Ma, \emph{Hilbert's sixth problem:
derivation of fluid equations via Boltzmann's kinetic theory},
arXiv:2503.01800 (preprint).
\bibitem{LazaroviciPickl} D.~Lazarovici, P.~Pickl, \emph{A mean field limit
for the Vlasov--Poisson system}, Arch.\ Ration.\ Mech.\ Anal.\ \textbf{225}
(2017), no.~3.
\bibitem{DuerinckxJabin} M.~Duerinckx, P.-E. Jabin, \emph{Derivation of 2D
Vlasov--Poisson for classical particles with Coulomb interactions},
arXiv:2608.04104 (2026, preprint).
\bibitem{DuerinckxSR} M.~Duerinckx, L.~Saint-Raymond, \emph{Lenard--Balescu
correction to mean-field theory}, Probab.\ Math.\ Phys.\ \textbf{2} (2021)
27--69.
\bibitem{DuerinckxLeBihan} M.~Duerinckx, C.~Le~Bihan, \emph{Lenard--Balescu
thermalization: rigorous derivation from a toy model}, arXiv:2511.10778
(2025, preprint).
\bibitem{GuoJang} Y.~Guo, J.~Jang, \emph{Global Hilbert expansion for the
Vlasov--Poisson--Boltzmann system}, Comm.\ Math.\ Phys.\ \textbf{299} (2010)
469--501.
\bibitem{BardosBesse} C.~Bardos, N.~Besse, \emph{Diffusion limit of the
Vlasov equation in the weak turbulent regime}, J.\ Math.\ Phys.\ \textbf{62}
(2021) 101505.
\bibitem{SHQ03} P.~Sharma, G.~W. Hammett, E.~Quataert, \emph{Transition from
collisionless to collisional MRI}, Astrophys.\ J.\ \textbf{596} (2003)
1121--1130.
\bibitem{SHD97} P.~B. Snyder, G.~W. Hammett, W.~Dorland, \emph{Landau fluid
models of collisionless magnetohydrodynamics}, Phys.\ Plasmas \textbf{4}
(1997) 3974.
\bibitem{JosephDimits} I.~Joseph, A.~M. Dimits, \emph{Connecting
collisionless Landau fluid closures to collisional plasma physics models},
Contrib.\ Plasma Phys.\ \textbf{56} (2016) 504 (abstract only consulted).
\bibitem{JiHeld} J.-Y. Ji, E.~D. Held, \emph{Electron parallel closures for
arbitrary collisionality}, Phys.\ Plasmas \textbf{21} (2014) 122116.
\bibitem{KK} F.~Kogelbauer, I.~Karlin, \emph{Exact hydrodynamic manifolds for
the linear Boltzmann BGK equation I: spectral theory}, Continuum Mech.\
Thermodyn.\ (2024), arXiv:2301.03069; and \emph{Spectral closure for the
linear Boltzmann-BGK equation}, arXiv:2306.07103.
\bibitem{Brodrick} J.~P. Brodrick et al., \emph{Testing nonlocal models of
electron thermal conduction for magnetic and inertial confinement fusion
applications}, arXiv:1704.08963 (2017).
\end{thebibliography}

\end{document}
